\section{Contesto}\label{Contesto}
Nei primi anni del 2000 il mondo stava vivendo in pieno la crescita esponenziale dell’ Internet che, un tempo limitato a università e ambiti militari, stava rapidamente entrando nelle case e in tutti gli uffici aziendali. Microsoft Windows e il pacchetto Office detenevano il monopolio dei servizi informatici. La posta elettronica, gestita prevalentemente tramite client come Microsoft Outlook, era diventata lo strumento di comunicazione globale per eccellenza. Di conseguenza un malware progettato per questi ambienti riusciva a replicarsi ovunque. 
Questa rapidissima digitalizzazione, infatti, non era stata accompagnata da un'adeguata consapevolezza dei rischi. Gli utenti si fidavano ciecamente dei messaggi ricevuti, e i sistemi operativi dell'epoca, come Windows 98, non possedevano le barriere di protezione odierne, arrivando persino a nascondere le vere estensioni dei file per impostazione predefinita.\\
Era il 4 maggio 2000 quando su una buona parte di computer di tutto il mondo arrivò una email con oggetto una dichiarazione d'amore "ILOVEYOU".
Dietro di essa si nascondeva un virus che infettò 45 milioni di computer in sole 24 ore. Il responsabile dell'attacco era uno studente di 23 anni, Onel de Guzman.
Egli, cresciuto in un contesto di povertà a Manila,nelle Filippine, riteneva che Internet dovesse essere accessibile a tutti. A questo proposito, per la sua tesi di laurea all' AMA Computer College, aveva lavorato ad un software che permetteva di ottenere password di accesso gratuito alla rete, ma il suo progetto fu bocciato. Così decise di modificarlo, aggiungendo un meccanismo di autopropagazione. Nella notte del 4 maggio lo inviò a Singapore, ignaro del fatto che l'indomani il mondo si sarebbe svegliato nel caos. Il panico non coinvolse solo privati, ma persino Microsoft, il Pentagono, il Parlamento Britannico e la NASA furono costretti a disconnettere i propri sistemi per arginare il contagio. Questi attacchi erano così innovativi che nelle FIlippine non esistevano delle leggi contro la  cybercriminalità e De Guzman non ebbe ripercussioni legali. 
Questo non fu il primo attacco informatico e nemmeno il più complesso dell’epoca, ma di certo dimostrò come il fattore umano potesse giocare un ruolo fondamentale nella sicurezza informatica.
